% GENERATED FILE. Edit canonical lessons, then run npm run generate:books.
% canonical-source-hash: fnv1a64:5271066d42056fc3
% canonical-lessons: GU-C495-sagarbha, GU-C495-prasuti, GU-C495-simant, GU-C495-avtarvu, GU-C495-janvu

\chapter{A Baby on the Way}
\label{ch:gu-a-baby-on-the-way}
\hlchaptermodality{495}

\begin{chapteropening}
\textbf{By the end of this chapter:} \emph{I can talk about a pregnancy and a birth.}

Complete the last lesson of chapter 495: I can talk about a pregnancy and a birth.
\end{chapteropening}

\section[\texorpdfstring{\gu{સગર્ભા} (sagarbhā)}{સગર્ભા (sagarbhā)}]{\gu{સગર્ભા} (sagarbhā) — pregnant; a pregnant woman}

\label{lesson:GU-C495-sagarbha}



\begin{quote}
Before the new one: say the Gujarati for elderly, then the Gujarati for an old people's home.
\end{quote}

\subsection*{You'll want to know: \gu{સગર્ભા}}
\textbf{\gu{સગર્ભા}} — \emph{sagarbhā} — \textquotedblleft{}pregnant; a pregnant woman\textquotedblright{}.

\begin{quote}
\gu{મારી} \gu{ભાભી} \gu{સગર્ભા} \gu{છે}.
\end{quote}

\begin{quote}
\emph{mārī bhābhī sagarbhā che.}
\end{quote}

\begin{quote}
My sister-in-law is pregnant.
\end{quote}

A feminine noun.

\begin{cousinweb}
\textbf{\gu{સગર્ભા}} (\emph{sagarbhā}) is a tatsama, Sanskrit \emph{sagarbhā}, \textquotedblleft{}a pregnant woman\textquotedblright{}, as Monier-Williams records it: \emph{sa-}, \textquotedblleft{}with\textquotedblright{}, and \emph{garbha}, \textquotedblleft{}womb, child\textquotedblright{}.
\end{cousinweb}

\begin{grammarlens}[title={What you've built}]
One more word for describing what you lived and felt.
\end{grammarlens}

\subsection*{Your turn}
\begin{itemize}
\raggedright
  \item \emph{Say it:} \emph{sagarbhā}
  \item \emph{Say it:} \emph{sagarbhā}, once more
  \item \emph{Say it:} \emph{mārī bhābhī sagarbhā che.} — My sister-in-law is pregnant.
  \item \emph{Recall:} say the Gujarati for elderly, then the Gujarati for an old people's home, then say \emph{sagarbhā} again
\end{itemize}

\begin{tcolorbox}[breakable,colback=teal!4,colframe=teal!35!black,title={Before you move on}]
What does \gu{સગર્ભા} mean? (\textquotedblleft{}Pregnant; a pregnant woman\textquotedblright{}.) Say it once more.
\end{tcolorbox}
\section[\texorpdfstring{\gu{પ્રસૂતિ} (prasūti)}{પ્રસૂતિ (prasūti)}]{\gu{પ્રસૂતિ} (prasūti) — childbirth; a delivery}

\label{lesson:GU-C495-prasuti}



\begin{quote}
Before the new one: say the Gujarati for an old people's home, then the Gujarati for pregnant.
\end{quote}

\subsection*{You'll want to know: \gu{પ્રસૂતિ}}
\textbf{\gu{પ્રસૂતિ}} — \emph{prasūti} — \textquotedblleft{}childbirth; a delivery\textquotedblright{}.

\begin{quote}
\gu{ભાભીની} \gu{પ્રસૂતિ} \gu{હોસ્પિટલમાં} \gu{થઈ}.
\end{quote}

\begin{quote}
\emph{bhābhīnī prasūti hospiṭalmā̃ thaī.}
\end{quote}

\begin{quote}
My sister-in-law gave birth in hospital.
\end{quote}

A feminine noun.

\begin{cousinweb}
\textbf{\gu{પ્રસૂતિ}} (\emph{prasūti}) is a tatsama, Sanskrit \emph{prasūti}, \textquotedblleft{}procreation, birth\textquotedblright{}, as Monier-Williams records it.
\end{cousinweb}

\begin{grammarlens}[title={What you've built}]
One more word for describing what you lived and felt.
\end{grammarlens}

\subsection*{Your turn}
\begin{itemize}
\raggedright
  \item \emph{Say it:} \emph{prasūti}
  \item \emph{Say it:} \emph{prasūti}, once more
  \item \emph{Say it:} \emph{bhābhīnī prasūti hospiṭalmā̃ thaī.} — My sister-in-law gave birth in hospital.
  \item \emph{Recall:} say the Gujarati for an old people's home, then the Gujarati for pregnant, then say \emph{prasūti} again
\end{itemize}

\begin{tcolorbox}[breakable,colback=teal!4,colframe=teal!35!black,title={Before you move on}]
What does \gu{પ્રસૂતિ} mean? (\textquotedblleft{}Childbirth; a delivery\textquotedblright{}.) Say it once more.
\end{tcolorbox}
\section[\texorpdfstring{\gu{સીમંત} (sīmant)}{સીમંત (sīmant)}]{\gu{સીમંત} (sīmant) — a baby-shower ceremony}

\label{lesson:GU-C495-simant}



\begin{quote}
Before the new one: say the Gujarati for pregnant, then the Gujarati for childbirth.
\end{quote}

\subsection*{You'll want to know: \gu{સીમંત}}
\textbf{\gu{સીમંત}} — \emph{sīmant} — \textquotedblleft{}a baby-shower ceremony\textquotedblright{}.

\begin{quote}
\gu{સાતમા} \gu{મહિને} \gu{ભાભીનું} \gu{સીમંત} \gu{થયું}.
\end{quote}

\begin{quote}
\emph{sātmā mahine bhābhīnũ sīmant thayũ.}
\end{quote}

\begin{quote}
In her seventh month my sister-in-law had her baby-shower ceremony.
\end{quote}

A neuter noun.

\begin{cousinweb}
\textbf{\gu{સીમંત}} (\emph{sīmant}) is a tatsama: Monier-Williams gives Sanskrit \emph{sīmanta} as \textquotedblleft{}the parting of the hair\textquotedblright{} and also as the hair-parting rite held for a woman expecting a child; Gujarati uses the word for that ceremony.
\end{cousinweb}

\begin{grammarlens}[title={What you've built}]
One more word for describing what you lived and felt.
\end{grammarlens}

\subsection*{Your turn}
\begin{itemize}
\raggedright
  \item \emph{Say it:} \emph{sīmant}
  \item \emph{Say it:} \emph{sīmant}, once more
  \item \emph{Say it:} \emph{sātmā mahine bhābhīnũ sīmant thayũ.} — In her seventh month my sister-in-law had her baby-shower ceremony.
  \item \emph{Recall:} say the Gujarati for pregnant, then the Gujarati for childbirth, then say \emph{sīmant} again
\end{itemize}

\begin{tcolorbox}[breakable,colback=teal!4,colframe=teal!35!black,title={Before you move on}]
What does \gu{સીમંત} mean? (\textquotedblleft{}A baby-shower ceremony\textquotedblright{}.) Say it once more.
\end{tcolorbox}
\section[\texorpdfstring{\gu{અવતરવું} (avtarvũ)}{અવતરવું (avtarvũ)}]{\gu{અવતરવું} (avtarvũ) — to come into the world; to be born}

\label{lesson:GU-C495-avtarvu}



\begin{quote}
Before the new one: say the Gujarati for childbirth, then the Gujarati for a baby-shower ceremony.
\end{quote}

\subsection*{You'll want to know: \gu{અવતરવું}}
\textbf{\gu{અવતરવું}} — \emph{avtarvũ} — \textquotedblleft{}to come into the world; to be born\textquotedblright{}.

\begin{quote}
\gu{ગઈકાલે} \gu{રાતે} \gu{અમારા} \gu{ઘરે} \gu{દીકરી} \gu{અવતરી}.
\end{quote}

\begin{quote}
\emph{gaīkāle rāte amārā ghare dīkrī avtarī.}
\end{quote}

\begin{quote}
Last night a baby daughter was born in our home.
\end{quote}

Past (it …): \gu{તે} \gu{બાળક} \gu{અવતર્યું}, \emph{te bāḷak avtaryũ}.

\begin{cousinweb}
\textbf{\gu{અવતરવું}} (\emph{avtarvũ}) goes back to Sanskrit \emph{ava-tṛ}, \textquotedblleft{}to come down, to descend\textquotedblright{} (Monier-Williams), the verb behind \emph{avatāra}: a child who comes down into the world.
\end{cousinweb}

\begin{grammarlens}[title={What you've built}]
One more word for describing what you lived and felt.
\end{grammarlens}

\subsection*{Your turn}
\begin{itemize}
\raggedright
  \item \emph{Say it:} \emph{avtarvũ}
  \item \emph{Say it:} \emph{avtarvũ}, once more
  \item \emph{Say it:} \emph{gaīkāle rāte amārā ghare dīkrī avtarī.} — Last night a baby daughter was born in our home.
  \item \emph{Recall:} say the Gujarati for childbirth, then the Gujarati for a baby-shower ceremony, then say \emph{avtarvũ} again
\end{itemize}

\begin{tcolorbox}[breakable,colback=teal!4,colframe=teal!35!black,title={Before you move on}]
What does \gu{અવતરવું} mean? (\textquotedblleft{}To come into the world; to be born\textquotedblright{}.) Say it once more.
\end{tcolorbox}
\section[\texorpdfstring{\gu{જણવું} (jaṇvũ)}{જણવું (jaṇvũ)}]{\gu{જણવું} (jaṇvũ) — to give birth to}

\label{lesson:GU-C495-janvu}



\begin{quote}
Before the new one: say the Gujarati for a baby-shower ceremony, then the Gujarati for to come into the world.
\end{quote}

\subsection*{You'll want to know: \gu{જણવું}}
\textbf{\gu{જણવું}} — \emph{jaṇvũ} — \textquotedblleft{}to give birth to\textquotedblright{}.

\begin{quote}
\gu{ગાયે} \gu{નાનું} \gu{વાછરડું} \gu{જણ્યું}.
\end{quote}

\begin{quote}
\emph{gāye nānũ vācharḍũ jaṇyũ.}
\end{quote}

\begin{quote}
The cow gave birth to a little calf.
\end{quote}

Past, with \gu{તેણે} (he or she …): \gu{તેણે} \gu{દીકરો} \gu{જણ્યો}, \emph{teṇe dīkro jaṇyo}.

\begin{cousinweb}
\textbf{\gu{જણવું}} (\emph{jaṇvũ}) is inherited from Sanskrit \emph{janayati}, \textquotedblleft{}causes to be born, begets\textquotedblright{}; CDIAL lists the Gujarati verb, \textquotedblleft{}to bear\textquotedblright{} (CDIAL 5102).
\end{cousinweb}

\begin{grammarlens}[title={What you've built}]
One more word for describing what you lived and felt.
\end{grammarlens}

\subsection*{Your turn}
\begin{itemize}
\raggedright
  \item \emph{Say it:} \emph{jaṇvũ}
  \item \emph{Say it:} \emph{jaṇvũ}, once more
  \item \emph{Say it:} \emph{gāye nānũ vācharḍũ jaṇyũ.} — The cow gave birth to a little calf.
  \item \emph{Recall:} say the Gujarati for a baby-shower ceremony, then the Gujarati for to come into the world, then say \emph{jaṇvũ} again
\end{itemize}

\begin{tcolorbox}[breakable,colback=teal!4,colframe=teal!35!black,title={Before you move on}]
What does \gu{જણવું} mean? (\textquotedblleft{}To give birth to\textquotedblright{}.) Say it once more.
\end{tcolorbox}
